% ============================================================
% Proposition 3.7: Error Propagation
% Status: FULL RECONSTRUCTION
% ============================================================

% --- Definitions and Assumptions ---
\begin{definition}[Estimation Error]
Let $\delta_\sigma(d,t) = \tilde{\sigma}_A(d,t) - \hat{\sigma}_A(d,t)$
be the Stage~1--Stage~2 discrepancy (Proposition~3.6). The true schema
coherence is $\sigma_A^*(d,t)$; the operative value used in the ODE is
$\tilde{\sigma}_A(d,t)$.
\end{definition}

\begin{definition}[Propagated ODE Error]
Let $\dot{\delta}_A^{\text{error}}(d,t)$ be the deviation in the depth
and schema ODE right-hand sides caused by using $\tilde{\sigma}_A$
instead of $\sigma_A^*$:
\[
\dot{\delta}_A^{\text{error}} =
|f(\mathbf{x}; \tilde{\sigma}_A) - f(\mathbf{x}; \sigma_A^*)|,
\]
where $f$ denotes the vector field of the coupled ODE system.
\end{definition}

\begin{assumption}[Lipschitz Continuity of Coupling Terms]
The coupling functions $\sigma_A \mapsto \lambda_c \gamma_\sigma \sigma_A
\delta_A$ and $\sigma_A \mapsto \rho P_A \alpha_A \sigma_A$ are
Lipschitz continuous with constants $L_1 = \lambda_c \gamma_\sigma
\delta_A$ and $L_2 = \rho P_A \alpha_A$, respectively.
\end{assumption}

\begin{assumption}[Bounded Variables]
All state variables are bounded as established in Proposition~3.2:
$\sigma_A \in [0,1]$, $\delta_A \in [0, \Delta_{\max}]$,
$\alpha_A \in [0,1]$, $P_A \in [0, P_{\max}]$.
\end{assumption}

% --- Proposition Statement ---
\begin{proposition}[Error Propagation in ODE System]
\label{prop:error-prop}
Under Assumptions~1--2, the effective error propagated into the depth
ODE \eqref{eq:depth-ode} and schema ODE \eqref{eq:sigma-ode} is
bounded by:
\[
|\dot{\delta}_A^{\text{error}}| \leq
\lambda_c \cdot \gamma_\sigma \cdot |\delta_\sigma| \cdot \delta_A
+ \rho \cdot P_A \cdot \alpha_A \cdot |\delta_\sigma|.
\]
If $|\delta_\sigma| < 0.10$, the propagated error is negligible
relative to typical ODE step sizes ($\eta_{\max} \sim O(1)$).
\end{proposition}

% --- Proof ---
\begin{proof}
We trace the two pathways through which $\sigma_A$ enters the
coupled ODE system.

\paragraph{Step 1: Depth ODE perturbation.}
The depth ODE \eqref{eq:depth-ode} is:
\[
\dot{\delta}_A = f_{\text{learn}} \cdot \eta \cdot T_A
- \lambda_c \cdot (1 - \gamma_\sigma \cdot \sigma_A) \cdot \delta_A \cdot (1 - r_A).
\]
The only $\sigma_A$-dependent term is the decay component
$-\lambda_c (1 - \gamma_\sigma \sigma_A) \delta_A (1 - r_A)$. Let
$\tilde{\sigma}_A = \sigma_A^* + \delta_\sigma$ be the corrupted value.
The deviation in the depth ODE is:
\begin{align*}
\Delta_{\delta} &= |
[-\lambda_c (1 - \gamma_\sigma \tilde{\sigma}_A) \delta_A (1 - r_A)]
- [-\lambda_c (1 - \gamma_\sigma \sigma_A^*) \delta_A (1 - r_A)] | \\
&= \lambda_c \delta_A (1 - r_A) \cdot |\gamma_\sigma \tilde{\sigma}_A - \gamma_\sigma \sigma_A^*| \\
&= \lambda_c \gamma_\sigma \delta_A (1 - r_A) \cdot |\delta_\sigma|.
\end{align*}
Since $(1 - r_A) \in [0, 1]$, we have:
\[
\Delta_{\delta} \leq \lambda_c \gamma_\sigma \delta_A |\delta_\sigma|.
\]
This is the first term in the bound.

\paragraph{Step 2: Schema ODE perturbation.}
The schema ODE \eqref{eq:sigma-ode} is:
\[
\dot{\sigma}_A = \sigma_A \cdot
[\rho P_A \alpha_A (1 - \gamma_\sigma \sigma_A) - \epsilon_\sigma \Omega_{SL}].
\]
The $\sigma_A$ factor outside the bracket is the operative value
(the state variable itself, not the estimate). The error propagation
occurs through the $\sigma_A$ inside the bracket — specifically, the
$(1 - \gamma_\sigma \sigma_A)$ which should use $\sigma_A^*$.

Let $\tilde{\sigma}_A = \sigma_A^* + \delta_\sigma$ be used in the
bracket's $(1 - \gamma_\sigma \sigma_A)$ term. The deviation in the
schema ODE at a given state is:
\begin{align*}
\Delta_{\sigma} &= |
\sigma_A [\rho P_A \alpha_A (1 - \gamma_\sigma \tilde{\sigma}_A) - \epsilon_\sigma \Omega_{SL}]
- \sigma_A [\rho P_A \alpha_A (1 - \gamma_\sigma \sigma_A^*) - \epsilon_\sigma \Omega_{SL}] | \\
&= \sigma_A \cdot \rho P_A \alpha_A \cdot |\gamma_\sigma \tilde{\sigma}_A - \gamma_\sigma \sigma_A^*| \\
&= \sigma_A \cdot \rho P_A \alpha_A \gamma_\sigma \cdot |\delta_\sigma|.
\end{align*}
Using $\sigma_A \leq 1$ and $\gamma_\sigma \leq 1$:
\[
\Delta_{\sigma} \leq \rho P_A \alpha_A |\delta_\sigma|.
\]
This is the second term in the bound.

\paragraph{Step 3: Combining terms.}
The total propagated error magnitude is the sum of the two pathways:
\[
|\dot{\delta}_A^{\text{error}}| = \Delta_{\delta} + \Delta_{\sigma}
\leq \lambda_c \gamma_\sigma \delta_A |\delta_\sigma|
+ \rho P_A \alpha_A |\delta_\sigma|,
\]
which is the stated inequality.

\paragraph{Step 4: Negligibility condition.}
The typical ODE step size is determined by the fastest rate constant
$\eta_{\max} \sim O(1)$. The propagated error is negligible when:
\[
\lambda_c \gamma_\sigma \delta_A |\delta_\sigma|
+ \rho P_A \alpha_A |\delta_\sigma| \ll \eta_{\max}.
\]
Using the numerical values $\lambda_c = 0.01$, $\gamma_\sigma \in (0,1)$,
$\rho = 0.1$, $\alpha_A \leq 1$, $P_A \leq 1$, $\delta_A \leq 1$:
\[
\lambda_c \gamma_\sigma \delta_A |\delta_\sigma| \leq 0.01 |\delta_\sigma|,
\qquad
\rho P_A \alpha_A |\delta_\sigma| \leq 0.1 |\delta_\sigma|.
\]
The total is $\leq 0.11 |\delta_\sigma|$. For $|\delta_\sigma| < 0.10$:
\[
|\dot{\delta}_A^{\text{error}}| \leq 0.011 \ll 1 \approx \eta_{\max}.
\]
Thus the propagated error is at most $1.1\%$ of the fastest rate,
rendering it negligible for numerical integration and qualitative
dynamics. $\square$
\end{proof}

% --- Boundary Cases ---
\begin{boundary-case}
\textbf{$\delta_A = 0$ (nascent domain):}
The depth ODE perturbation term vanishes ($\Delta_{\delta} = 0$).
Only the schema ODE perturbation remains, bounded by
$\rho P_A \alpha_A |\delta_\sigma|$.

\textbf{$\alpha_A = 0$ (attention collapse):}
Both perturbation terms vanish. Errors in $\sigma_A$ estimation
do not propagate when there is no attentional gating. This is
consistent: if $\alpha_A = 0$, the ODE dynamics are already
frozen.

\textbf{$P_A = 0$ (no prior knowledge):}
The schema perturbation term vanishes. Only depth perturbation,
bounded by $\lambda_c \gamma_\sigma \delta_A |\delta_\sigma|$, remains.

\textbf{$|\delta_\sigma| \geq 0.10$:}
The propagated error may become non-negligible. The recalibration
rule (ridge regression, Proposition~3.6) is triggered to reduce
$|\delta_\sigma|$ below this threshold, restoring negligibility.
\end{boundary-case}

% --- Circular Dependency Check ---
\begin{circularity-check}
\textbf{Dependencies on earlier results:} The proof uses
Proposition~3.6 to define $\delta_\sigma$ and its bound. It also
uses the ODE definitions from Section~3.2 and Proposition~3.2
(boundedness of state variables). These are all earlier results.

\textbf{Forward dependencies:} The error bound is used implicitly
in Theorem~4.1 to ensure that SGD suppression analysis is robust
to estimation noise. This is a forward dependency, not circular.

\textbf{Self-circularity:} The bound assumes the error propagates
additively through the two pathways. In reality, the ODE is coupled
and errors can amplify through feedback — $\delta_\sigma$ affects
$\dot{\sigma}_A$, which changes $\sigma_A(t + \Delta t)$, which
changes $\delta_\sigma$ at the next timestep. The stated bound is a
one-step bound (instantaneous), not a global bound over a time
interval. A full global error bound would require a Gronwall-type
argument integrating the Lipschitz constant over time. The
one-step bound is sufficient for the claim's purpose (assessing
negligibility relative to $\eta_{\max}$).
\end{circularity-check}
